\clearpage
\VerdierPhysicalPageAnchor{0233}%
\begingroup\small\setlength{\parskip}{0pt}\setlength{\emergencystretch}{1.5em}
\setlength{\abovedisplayskip}{4pt plus 2pt minus 2pt}%
\setlength{\belowdisplayskip}{4pt plus 2pt minus 2pt}%

\VerdierSemanticAnchor{ch4d-4-1-8}
\noindent\textbf{4.1.8.}\quad
Lorsque $X$ est un complexe simple, l'objet spectral à valeurs dans
$K(\mathcal A)$ défini par le complexe $X$ filtré par son premier degré
(4.1.7) n'est autre que le premier objet spectral à valeurs dans
$K(\mathcal A)$ associé au complexe $X$.

\VerdierSemanticAnchor{ch4d-4-1-8-degree-support}
De même, soient $X$ un complexe $n$-uple et $m$ un entier,
$1\leq m\leq n$. Supposons que le composant $X^\sigma$ de $X$ soit nul
lorsque :
\[
\sigma-\sigma_m e_m\neq0.
\]
Ici, $\sigma_m$ désigne la $m$-ième coordonnée.
On vérifie alors que l'objet spectral défini par le complexe $X$ filtré par
son $m$-ième degré est égal au premier objet spectral canonique du complexe
simple $\int_m X$.

\VerdierSemanticAnchor{ch4d-4-1-8-derived-image}
Enfin, lorsque la catégorie $\mathcal A$ est abélienne, le foncteur
canonique :
\[
K(\mathcal A)\longrightarrow D(\mathcal A)
\]
transforme l'objet spectral à valeurs dans $K(\mathcal A)$ défini par un
complexe $n$-uple $X$ filtré par son $m$-ième degré en un objet spectral à
valeurs dans $D(\mathcal A)$. L'objet spectral ainsi obtenu est appelé
\emph{l'objet spectral, à valeurs dans $D(\mathcal A)$, défini par $X$,
filtré par son $m$-ième degré}. Lorsqu'aucune confusion n'en résulte, on
supprime la mention ``à valeurs dans $D(\mathcal A)$''.

\VerdierSemanticAnchor{ch4d-4-1-9}
\medskip\noindent\textbf{4.1.9.}\quad
Soit $X^{\bullet\bullet}$ un complexe double de $\mathcal A$.

\VerdierSemanticAnchor{ch4d-4-1-9-row-complexes}
Désignons, pour tout entier $q$, par $X^{q,\bullet}$ (resp.
$X^{\bullet,q}$) le complexe simple :
\[
\begin{aligned}
X^{q,\bullet}
&=\cdots\xrightarrow{d_2^{q,i-1}}X^{q,i}
  \xrightarrow{d_2^{q,i}}X^{q,i+1}
  \xrightarrow{d_2^{q,i+1}}\cdots,\\[3pt]
(\text{resp.}\quad X^{\bullet,q}
&=\cdots\longrightarrow X^{i,q}
  \xrightarrow{(-1)^q d_1^{i,q}}X^{i+1,q}
  \xrightarrow{(-1)^q d_1^{i+1,q}}X^{i+2,q}
  \longrightarrow\cdots\quad),
\end{aligned}
\]
dont le composant de degré $i$ est $X^{q,i}$ (resp. $X^{i,q}$).

\VerdierSemanticAnchor{ch4d-4-1-9-direction-morphisms}
La différentielle dans la première (resp.\ deuxième) direction de
$X^{\bullet\bullet}$ induit des morphismes (dans $K(\mathcal A)$) :
\[
\begin{aligned}
d_1^{q,\bullet}&\colon X^{q,\bullet}\longrightarrow X^{q+1,\bullet}\\
(\text{resp.}\quad
d_2^{\bullet,q}&\colon X^{\bullet,q}\longrightarrow X^{\bullet,q+1}\quad),
\end{aligned}
\]

\VerdierSemanticAnchor{ch4d-4-1-9-components}
dont les composants sont :
\[
\begin{aligned}
(d_1^{q,\bullet})^i&=d_1^{q,i}\\
(\text{resp.}\quad(d_2^{\bullet,q})^i&=(-1)^i d_2^{i,q}\quad).
\end{aligned}
\]
\endgroup
